\documentclass[10pt,a4paper]{amsart}
\usepackage[margin=19mm]{geometry}
\usepackage{amsmath,amssymb,mathtools}
\usepackage[scr=rsfs,cal=euler]{mathalfa}
\usepackage{xfrac}
\usepackage[unicode,hidelinks,bookmarks=false]{hyperref}
\newcommand{\ie}{i.e.\ }
\newcommand{\cf}{cf.\ }
\newcommand{\resp}{respectively\ }
\newcommand{\Subsection}{\subsection}
\providecommand{\nobreakdash}{\nobreak\hbox{-}}
\setlength{\emergencystretch}{2em}
\multlinegap0pt
\usepackage{amssymb}
\usepackage{mathtools}
\usepackage{xcolor}
\newtheorem{proposition}{Proposition}
\title{Entropy bounds and global couplings\\for union-closed families}
\author{Yunjiang Jiang}
\date{Source: arXiv:2609.08291v1 (CC BY 4.0)}
\begin{document}
\maketitle

\noindent\emph{Source and licence.} Entropy bounds and global couplings\\for union-closed families. Version arXiv:2609.08291v1 (CC BY 4.0), distributed by arXiv under the Creative Commons Attribution 4.0 International licence, http://creativecommons.org/licenses/by/4.0/, as recorded in the arXiv licence metadata for this exact version and shown on the abstract page https://arxiv.org/abs/2609.08291v1. Full licence notice: \texttt{licenses/arxiv-2609.08291-CC-BY-4.0.txt}.\par\smallskip

\noindent\textit{Editorial scope.} Counted unit: the article's proposition prop:sharper-ceiling (source0.tex lines 414--425). The article's complete original proof of it is delivered with it (source0.tex lines 426--498, one \texttt{proof} environment of 73 lines). No mathematics is written, repaired or re-derived: the statement and the proof are the article's own lines, in the article's own notation, and no retained line is substituted or re-flowed.  Retained uncounted as the source-local context the proof needs: lines 147--149; lines 151--154; lines 280--284.  Nothing was omitted from any retained span, so the package carries no omission record.  No retained line contains a citation, so the wrapper carries no bibliography and no bibliography entry.  No retained line contains a figure environment, an image-inclusion command or drawing code, so no image is delivered and the package declares no asset dependency.  Editorial formatting only, in addition: an AMS \texttt{amsart} preamble that restates the article's own declarations and macros used by the retained lines, namely the environments \texttt{proposition} on separate counters, together with the standard packages those lines need.  The retained environments print in this document as Proposition 1 in order of appearance, whereas the article prints the same items with the numbering of its own full document.  The complete native source of the article as distributed in the arXiv e-print for this exact version is bundled unchanged as \texttt{sources/209707/source0.tex}.
\begin{equation}\label{eq:stationary}
(1-q)\ln(1-q)+q\ln(1+q)=0.
\end{equation}

\begin{equation}\label{eq:constant-ceiling}
C_*=\frac12-\frac{(1-q_*)\ln(1+q_*)}{2\ln2}
\approx0.3828852599667978.
\end{equation}

\begin{align}
K(x,y)&\ge\psi(x)\psi(y),\tag{K}\label{eq:K}\\
2m\bigl[w_0h(xy)+w_4\psi(x)\psi(y)\bigr]
&\ge xh(y)+yh(x).\tag{M}\label{eq:M}
\end{align}

\begin{proposition}[Limitation of the pointwise method]
\label{prop:sharper-ceiling}
If \eqref{eq:K}--\eqref{eq:M} hold, then
\[
1-m\le C_*:=
\frac12-\max_{0<q<1}
\frac{qh(q)-\frac12h(q^2)}{\ln2}.
\]
The maximum occurs at the unique solution of
\eqref{eq:stationary}, and \(C_*\) is given by
\eqref{eq:constant-ceiling}.
\end{proposition}

\begin{proof}
The sampling kernel satisfies \(K(1,1)=h(1)=0\).
Thus \eqref{eq:K} forces \(\psi(1)=0\).
Put \(y=1\) in \eqref{eq:M}, and choose any \(0<x<1\). It gives
\[
2mw_0h(x)\ge h(x),
\qquad\text{hence}\qquad mw_0\ge\frac12.
\tag{E}
\label{eq:endpoint-barrier}
\]
Next put \(x=y=q\) in \eqref{eq:M}. Since
\(\psi(q)^2\le K(q,q)\le\ln2\), and \(w_4=1-w_0\), we obtain
\[
\begin{aligned}
qh(q)
&\le m\bigl[w_0h(q^2)+w_4\psi(q)^2\bigr]\\
&\le m\ln2-mw_0\bigl[\ln2-h(q^2)\bigr]\\
&\le m\ln2-\frac12\bigl[\ln2-h(q^2)\bigr].
\end{aligned}
\]
The last step uses (E) and \(h(q^2)\le\ln2\).
Rearranging gives, for every \(q\in(0,1)\),
\[
1-m\le\frac12-\frac{qh(q)-\frac12h(q^2)}{\ln2}.
\tag{C}
\label{eq:ceiling-each-q}
\]
Taking the strongest of these bounds proves the claimed formula.

To identify the maximizing point, put \(f(q)=qh(q)-h(q^2)/2\).
Direct differentiation and cancellation give
\[
f'(q)=-(1-q)\ln(1-q)-q\ln(1+q),
\]
\[
f''(q)=\ln\frac{1-q}{1+q}+\frac1{1+q}.
\]
For \(0<q\le1/2\), the inequalities
\(-\ln(1-q)>q\) and \(\ln(1+q)<q\) show that
\(f'(q)>q(1-2q)\ge0\).
For \(q\ge1/2\), \(f''(1/2)=-\ln3+2/3<0\), and
\[
f'''(q)=-\frac2{1-q^2}-\frac1{(1+q)^2}<0.
\]
Thus \(f'\) is strictly decreasing on \([1/2,1)\).
It starts there at \(\frac12\ln(4/3)>0\) and tends to
\(-\ln2<0\). There is exactly one maximizing point \(q_*\).
The equation \(f'(q_*)=0\) is precisely
\eqref{eq:stationary}. Expanding \(h(q^2)\) gives
\[
f(q)=\frac12(1-q)^2\ln(1-q)
     +\frac12(1-q^2)\ln(1+q).
\]
At the stationary point,
\((1-q_*)\ln(1-q_*)=-q_*\ln(1+q_*)\), so this simplifies to
\[
f(q_*)=\frac12(1-q_*)\ln(1+q_*).
\]
Substitution proves \eqref{eq:constant-ceiling}.
For numerical evaluation, put
\(\Phi(q)=(1-q)\ln(1-q)+q\ln(1+q)\).
Ordinary Newton iteration is
\[
q_{j+1}=q_j-
\frac{\Phi(q_j)}
{\displaystyle\ln\frac{1+q_j}{1-q_j}-\frac1{1+q_j}}.
\]
Starting at \(q_0=0.69\) gives
\(q_*\approx0.6909077387254078\), and substitution into
\eqref{eq:constant-ceiling} gives the stated decimal value of \(C_*\).
No formal certification of these decimal approximations is needed:
the formulas for \(q_*\) and \(C_*\) are exact.
\end{proof}

\end{document}
